
% [X] add more sources to Simulationsprozess
\section{Der Simulationsprozess}
 
Unter Simulation wird die Nachahmung der Funktionsweise eines realen Systems verstanden. Im Kerne der Simulation steht ein Modell, Konstrukt, das die für die Anwendung wichtigste Merkmale des ausgewählten Systems darstellt. 

Sie finden Anwendung in unterschiedliche Bereiche der Technik und nehmen eine wichtige Rolle bei der Lösung dynamischen und komplexen Problemen, die mittels analytischer Vorgehensweise nicht ohne weiteres gelöst werden können. Durch ihren Einsatz gewinnt der Nutzer Klarheit und Verständnis an dem Entwicklungsprozess \cite{zimmerDurchgaengigerSimulationsprozessZur2015}. 

In den folgenden Unterkapiteln werden \textit{Hardware-in-the-loop} und \textit{Software-in-the-loop} vorgestellt, zwei Prozesse der Simulationstechnik, die im Produktentstehungsprozess häufig verwendet werden. Der Begriff \textit{x-in-the-loop} deutet darauf hin, dass Abstraktionen eines Systems im Laufe des Entwicklungsprozesses eingesetzt werden, bzw. die Anbindung einer Hardware, sowie eine Software-Darstellung des Systems an die Simulationsumgebung.

% RESEARCH
\begin{comment}
Simulationen werden häufig eingesetzt, wenn dynamische Vorgänge eines Systems beschrieben werden sollen, ein Lösungsansatz nach analytischer Vorgehensweise nicht ohne weiteres gefunden werden kann oder ein Lösungsansatz zu stark simplifiziert ist.
Der Begriff „Simulation“ wird oftmals mit den Wörtern „Verifikation“ und „Validierung“ in Verbindung gebracht. Gemeint ist, dass die einzelnen Simulationsmodelle der durchzuführenden Simulation den Anforderungen, welche man an sie stellt, auch entsprechen

Validierung Durch die Validierung soll die Anforderung an ein Produkt und eine Wiederholbarkeit des Ergebnisses sichergestellt werden. Die Validierung wird durchgeführt, indem das Modellverhalten mit dem Verhalten des realen Systems verglichen wird und die Differenzen evaluiert werden.

Verifikation Die Verifikation soll gewährleisten, dass ein Simulationsmodell seiner Spezifikation entspricht und damit zum Beispiel die Wertigkeit eines Softwareproduktes für den operativen Einsatz aufzeigt

Im Folgenden werden Prozessabläufe einer Simulation exemplarisch vorgestellt, welche unter anderem bei der Umsetzung des in dieser Arbeit entwickelten Simulationsprozesses als Modell dienen. Durch die sukzessive Weiterentwicklung von Simulationsmethoden können Lücken in der Prozesskette geschlossen und neue Felder identifiziert werden. \cite{zimmerDurchgaengigerSimulationsprozessZur2015}

x-in-the-Loop
software-in-the-loop
model-in-the-loop

In Abbildung 30 sind die drei Eckpfeiler einer durchgängigen Simulationskette dargestellt. Dabei wird die Informationsnotwendigkeit durch die (verfügbaren) Modelle, das Konzept der Simulation durch die Architektur und die Zusammensetzung der einzelnen Simulationskomponenten durch die Anforderungen repräsentiert. Mit diesen Anforderungen an den Umgang der Durchgängigkeit lassen sich in dieser Produktentwicklungsphase ein besseres Gesamtverständnis des Systems und eine Steigerung des Reifegrades zu früheren Zeitpunkten erreichen.

\img{10cm}{media/2022-04-02-18-31-51.png}{simulationsprozess}{Die drei Eckpfeiler eines durchgängigen Simulationsprozesses: Modelle, Architektur und Anforderungen}


\end{comment}

% [X] Ergänzung Hardware-in-the-loop
\subsection{\textit{Hardware-in-the-Loop (HiL)}}

Ein reales System durch Einsatz von physischen Prototypen ist oftmals wegen der mehreren Konfigurationsmöglichkeiten nicht möglich. Eine Alternative dazu stellt die Anbindung der Hardware an die Simulationsumgebung durch eine elektrische Schnittstelle \cite{zimmerDurchgaengigerSimulationsprozessZur2015}. Dieser Prozess wird als \textit{Hardware-in-the-Loop} (\textit{HiL}) genannt. \textquote[\cite{zimmerDurchgaengigerSimulationsprozessZur2015}]{\textit{Hardware-in-the-Loop (HiL)} Ansatzes bezeichnet oftmals den Einsatz eines realen Steuergerätes oder eines realen Bauteils in einem virtuellen Fahrversuch, ohne dabei ein Risiko für den Fahrer oder Schäden an Prototypen einzugehen.}

\img[ellisControlSystemDesign2012]{10cm}{media/2022-08-22-11-39-16.png}{hilreg}{Blockschaltbild eines Regelungssystems, in dem der Regler im HIL getest ist}

HIL ist seit Jahren in verschieden Branchen der Industrie, von Aerospace bis zu Automotiv, eingesetzt, in denen Systeme besonders komplex sind. Mit seinem Einsatzt werden mehreren Vorteilen erzielt, wie der Reduzierung der Kosten zum Testen und des mit Ausfall einhergehenden Risikos. Außerdem lässt sich damit, Komponenten früher im Produktlebenszyklus testen, da andere für das Testen nowendigen Komponente eventuell nur später in der Entwicklung verfügbar sein können. \cite{ellisControlSystemDesign2012}

Ein anderer wichtiger Vorteil dieses Ansatzes bezieht sich auf das Variantenmanagement. Varianten können mit Einsatz von HIL leichter erzeugt werden, unter Berücksichtigung, dass Anlagen oft variieren können, und die Prüfung von Mustern jeder Variante kann teuer sein. Mit Hilfe von HIL können daher Tests nur an einer Teilmenge der verfügbaren Varianten durchgeführt werden. \cite{ellisControlSystemDesign2012}    


% RESEARCH
\begin{comment}
Durch die zunehmende Anzahl an vernetzten Steuergeräten in Fahrzeugen nimmt die Bedeutung der Simulation zur Steigerung von Synergiepotenzialen und zur Vermeidung von Fehlfunktionen stetig zu. Eine Überprüfung durch reale Prototypen ist durch die Vielzahl an Konfigurationsmöglichkeiten nicht mehr möglich. Außerdem sind häufig Einsatzszenarien gefordert, zu denen die Hardwarekomponente noch nicht verfügbar ist. Das Konzept des Hardware-in-the-Loop (HiL) Ansatzes bezeichnet oftmals den Einsatz eines realen Steuergerätes oder eines realen Bauteils in einem virtuellen Fahrversuch, ohne dabei ein Risiko für den Fahrer oder Schäden an Prototypen einzugehen.
Im Fokus sind bei diesem Ansatz vor allem Steuergeräte wie Fahrwerksregelsysteme, Antriebssteuergeräte und Komfortsteuergeräte sowie Bauteile von Motor, Getriebe, etc. des Fahrzeugs. \cite{zimmerDurchgangigerSimulationsprozessZur2015}

HIL has been used for years in industries such as aerospace or automotive where plants are especially complex. HIL can augment testing of physical plants to provide many advantages such as: • Reduced cost to test The high cost to execute tests on complex machinery like aircraft sub-systems can justify substantial investment in HIL. Field testing can be reduced by using HIL as a complement. On their Web site, National Instruments presents one example where Process Automation

\end{comment}

\subsection{\textit{Software-in-the-Loop}}
Ein Nachteil des \textit{Hardware-in-the-Loop} wäre die relative Verzögerung, bis im Projekt die Hardware verfügbar wird \cite{zimmerDurchgaengigerSimulationsprozessZur2015}. Eine Alternative für das Testen in früheren Phasen ist die Erstellung einer Software-Darstellung des Systems, die die für die Phase notwendigen Anforderungen abdeckt. Damit lässt sich eine reine Software-Schnittstelle zwischen den Testkomponenten verwenden, im Gegensatz zu einer aufwendigeren elektrischen Schnittstelle, was für den \textit{Hardware-in-the-Loop} der Fall ist \cite{zimmerDurchgaengigerSimulationsprozessZur2015}. 


% RESEARCH
\begin{comment}
Der Nachteil der oben erwähnten HiL-Systeme, dass die Hardware teilweise relativ spät zur Verfügung steht, kann durch den \textit{Software-in-the-loop (SiL)} Ansatz kompensiert werden. Während der Implementierungsphase können Nutzer Regelungsstrategien anhand dieser SiL Systemen testen. Im Gegensatz zu HiL Systemen, bei denen eine elektrische Schnittstelle verwendet wird, sind die Schnittstellen zwischen den Testkomponenten auf Betriebssystem-Ebene, also reine Software-Schnittstellen. \cite{zimmerDurchgaengigerSimulationsprozessZur2015}
\end{comment}


% \subsection{Model-in-the-Loop}

% % RESEARCH
% \begin{comment}
% Das Funktionsmodell, welches getestet werden soll, ist ein Modell des Systems. Damit befindet es sich innerhalb der Regelschleife und umfangreiche Funktions157 Eigene Darstellung. Simulationsgrundlagen tests in einer frühen Entwicklungsphase sind möglich. Zusätzlich zu den Funktionstests können Modellabdeckungstests durchgeführt werden. In dieser Arbeit werden für die Bewertung von Fahrzeugkonzepten unterschiedliche Funktionsmodelle in das Simulationsmodell implementiert, so dass Anforderungen an die Komponenten im Kontext Gesamtsystem untersucht werden können. \cite{zimmerDurchgangigerSimulationsprozessZur2015}
% \end{comment}


